%% Système bielle-manivelle : transformation d'une rotation continue en translation alternative.
%% Manivelle 1 (vert) : rotation autour de O (pivot au bâti 0) — le point A décrit un cercle
%% de rayon OA (trajectoire en pointillés), repéré par l'angle theta depuis l'axe (O,x).
%% Bielle 2 (rouge) : mouvement plan.  Piston 3 (bleu) : translation alternative dans son guidage.
\documentclass[tikz,border=4pt]{standalone}
\usepackage{liaisons}
\usetikzlibrary{backgrounds}
\begin{document}
\begin{tikzpicture}[x=1cm,y=1cm,
    solide/.style={line width=1.4pt, line cap=round, line join=round},
    axe/.style={line width=0.5pt, dash pattern=on 6pt off 2pt on 1pt off 2pt},
    fleche/.style={-{Stealth[length=5pt]}, line width=0.8pt}]
\useasboundingbox (-2.2,-2.1) rectangle (6.2,1.6);   % boîte fixe (animation : cm → SVG)
\colorlet{S0}{black}\colorlet{S1}{solideE}\colorlet{S2}{solideA}\colorlet{S3}{solideB}
% --- Géométrie : OA = 1,2 ; theta = 55° ; AB = 2,9 ---
\coordinate (O) at (0,0);
\coordinate (A) at (55:1.2);
\coordinate (B) at (3.62,0);
% --- Trajectoire du point A (cercle) et angle theta ---
\draw[line width=0.6pt, dashed, draw=S1!70] (O) circle (1.2);
\draw[axe] (-1.6,0) -- (5.8,0);
\draw[line width=0.7pt, draw=S1!60] (0.62,0) arc (0:55:0.62);
\node[font=\small, text=S1, anchor=west, inner sep=1.5pt] at (27:0.72) {$\theta$};
% --- Guidage du piston (bâti) : deux traits hachurés ---
\foreach \x in {1.4,1.6,...,5.25} {\draw[line width=0.5pt] (\x,0.5) -- (\x-0.14,0.64)
   (\x,-0.5) -- (\x-0.14,-0.64);}
\draw[line width=0.8pt] (1.25,0.5) -- (5.3,0.5) (1.25,-0.5) -- (5.3,-0.5);   % guidage : couvre toute la course du piston
% --- Symboles de liaison ---
\pic[s1=S1, s2=S0, liens=false] at (O) {pivot bout};
\pic[s1=S2, s2=S1, liens=false] at (A) {pivot bout};
\pic[s1=S2, s2=S3, liens=false] at (B) {pivot bout};
\pic at (0,-1.7) {bati};
% --- Solides ---
\begin{scope}[on background layer]
  \draw[solide, S0] (O) -- (0,-1.7);                                  % liaison au bâti 0
  \draw[solide, S1] (O) -- (A);                                       % manivelle 1
  \draw[solide, S2] (A) -- (B);                                       % bielle 2
  \draw[solide, S3, fill=white] (3.02,-0.42) rectangle (4.22,0.42);   % piston 3
\end{scope}
% --- Mouvements ---
\draw[fleche, S1!60] ([shift=(60:0.81)]O) arc (60:142:0.81);
\node[font=\small, text=S1, anchor=east, inner sep=2pt] at (-0.7,0.42) {$\omega$};
\draw[{Stealth[length=5pt]}-{Stealth[length=5pt]}, line width=0.8pt, S3!60] (3.2,-0.95) -- (4.8,-0.95);
\node[font=\small, text=S3, anchor=north] at (4.0,-1.05) {translation alternative};
% --- Points et repères des solides ---
\node[font=\small, anchor=north east, inner sep=3pt] at (O) {O};
\node[font=\small, anchor=south east, inner sep=3pt] at (A) {A};
\node[font=\small, anchor=south, inner sep=4pt] at (B) {B};
\node[font=\small\bfseries, text=S0, anchor=west] at (0.28,-1.62) {0};
\node[font=\small, text=S1, anchor=east] at (-0.42,-1.45) {\textbf{1} manivelle};
\node[font=\small, text=S2, anchor=south] at (2.15,0.6) {\textbf{2} bielle};
\node[font=\small, text=S3, anchor=south] at (4.25,0.78) {\textbf{3} piston};
\node[font=\small, text=S1!70, anchor=east, align=right] at (-0.78,1.02) {trajectoire\\de A};
\repere{(1.35,-1.75)}{x}{y}{1}
\end{tikzpicture}
\end{document}
